\documentclass[border=2pt,tikz]{standalone}
% Shared TikZ/pgfplots style for every figure in the manuscript.
\usepackage[T1]{fontenc}
\usepackage{mathpazo}
\usepackage{amsmath,amssymb}
\usepackage{tikz}
\usepackage{pgfplots}
\usepackage{pgfplotstable}
\pgfplotsset{compat=1.18}
\usetikzlibrary{arrows.meta,positioning,calc,backgrounds,fit,shapes.geometric}
\usepgfplotslibrary{fillbetween}

% --- palette: colour-blind safe, legible in greyscale -------------------
\definecolor{cHyd}{HTML}{1B6CA8}   % hydride / primary blue
\definecolor{cPro}{HTML}{C1492E}   % proton  / primary red
\definecolor{cAcc}{HTML}{7B5AA6}   % accent  / purple
\definecolor{cGold}{HTML}{B8860B}  % accent  / gold
\definecolor{cGrey}{HTML}{6E6E6E}
\definecolor{cInk}{HTML}{1A1A1A}
\definecolor{cBandR}{HTML}{C1492E}
\definecolor{cBandG}{HTML}{4F7942}

% --- base axis ----------------------------------------------------------
\pgfplotsset{
  icig/.style={
    width=68mm, height=52mm,
    scale only axis,
    axis line style={cInk, line width=0.45pt},
    tick style={cInk, line width=0.4pt},
    tick align=outside, tick pos=left,
    label style={font=\small, color=cInk},
    tick label style={font=\footnotesize, color=cInk},
    title style={font=\small\bfseries, color=cInk, align=left},
    legend cell align=left,
    legend style={draw=none, fill=none, font=\scriptsize,
                  inner sep=1pt, row sep=0.5pt},
    grid=both,
    major grid style={cGrey!22, line width=0.28pt},
    minor grid style={cGrey!12, line width=0.2pt},
  },
  % marker shorthands
  ptA/.style={mark=*, mark size=1.7pt, thick,
              mark options={fill=cHyd, draw=white, line width=0.4pt}},
  ptB/.style={mark=*, mark size=1.7pt, thick,
              mark options={fill=cPro, draw=white, line width=0.4pt}},
  hollowB/.style={mark=o, mark size=1.7pt, thick,
              mark options={draw=cPro, line width=0.7pt}},
}

% panel tag, e.g. \panel{a}
\newcommand{\panel}[1]{\textbf{(\textsf{#1})}}

\usepgfplotslibrary{fillbetween}
\begin{document}
\begin{tikzpicture}
% Envelope figure, full text width.  Every curve is computed by
% analysis/export_fig_envelope.py.
% Generated by analysis/export_fig_envelope.py; do not edit.
\def\feFzero{-0.042086}
\def\feXthr{0.238146}
\def\fePeakLo{140}
\def\fePeakHi{470}
\def\feCross{417.0}
\def\feYadhX{1.964311}
\def\feYadhF{0.128726}
\def\feLightX{1.118415}
\def\feLightF{-0.261685}

\pgfplotsset{
  env/.style={icig, scale only axis,
    label style={font=\fontsize{8.0}{9.6}\selectfont, color=cInk},
    tick label style={font=\fontsize{7.2}{8.6}\selectfont, color=cInk},
    title style={font=\fontsize{8.6}{10.4}\selectfont\bfseries, color=cInk,
                 align=left}},
}
\tikzset{
  ann/.style={font=\fontsize{7.2}{8.6}\selectfont, color=cInk, align=left,
              inner sep=0pt},
  annc/.style={font=\fontsize{7.2}{8.6}\selectfont, color=cInk, align=center,
              inner sep=0pt},
}

% =====================================================================
% (A) the attainable region against two compatible sets
% =====================================================================
\begin{axis}[env, name=A, width=50mm, height=50mm,
  xmin=0, xmax=3.0, ymin=-0.46, ymax=0.55,
  xlabel={$X_{\mathrm H}=\ln x_{\mathrm H}$, intrinsic},
  ylabel={intrinsic offset $F$},
  xtick={0,0.5,1,1.5,2,2.5,3},
  xticklabels={$0$,$0.5$,$1$,$1.5$,$2$,$2.5$,$3$},
  ytick={-0.4,-0.2,0,0.2,0.4},
  yticklabels={$-0.4$,$-0.2$,$0$,$0.2$,$0.4$},
  clip=true,
  title={(\textit{\textbf{A}})\ \ what the model attains, and what\\
         \hspace*{1.2mm}two observations allow}]
  \addplot[name path=ENV, draw=none] table[x=xh,y=F] {data/fe_env.dat};
  \addplot[name path=RIG, draw=none] table[x=xh,y=rigid] {data/fe_env.dat};
  \addplot[cAcc!15] fill between[of=ENV and RIG];
  \addplot[cAcc, line width=0.9pt] table[x=xh,y=F] {data/fe_env.dat};
  \addplot[cAcc!70, line width=0.6pt, densely dotted] table[x=xh,y=rigid]
    {data/fe_env.dat};
  \draw[cInk, line width=0.6pt, densely dashed] (axis cs:0,0) -- (axis cs:3,0);
  \draw[cGrey, line width=0.4pt, densely dotted]
    (axis cs:\feXthr,-0.46) -- (axis cs:\feXthr,0.55);

  \addplot[cPro, line width=1.1pt] table[x=xh,y=F] {data/fe_obs_yadh.dat};
  \addplot[cHyd, line width=1.1pt, densely dashed] table[x=xh,y=F]
    {data/fe_obs_light.dat};
  \addplot[only marks, mark=o, mark size=2.0pt, cPro, line width=0.7pt]
    coordinates {(\feYadhX,\feYadhF)};
  \addplot[only marks, mark=o, mark size=2.0pt, cHyd, line width=0.7pt]
    coordinates {(\feLightX,\feLightF)};

  \node[ann, anchor=south west] at (axis cs:0.30,0.012) {$F=0$, semiclassical};
  \node[ann, color=cAcc!85!black, anchor=south east, text width=32mm,
        align=right] at (axis cs:2.95,-0.445) {\mbox{ground-channel model:}\\attainable offsets};
  \node[ann, color=cAcc!85!black, anchor=north east] at (axis cs:2.95,-0.058)
    {envelope, $\Fz$};
  \node[ann, color=cAcc!70!black, anchor=south, rotate=-47.1]
    at (axis cs:0.62,-0.170) {rigid, $w\to0$};
  \node[ann, color=cGrey, anchor=west] at (axis cs:\feXthr+0.04,0.50)
    {$x_{\mathrm H}=1.27$};
  \node[ann, color=cPro, anchor=east] at (axis cs:1.93,0.27) {yeast ADH};
  \node[ann, color=cHyd, anchor=north west] at (axis cs:1.17,-0.275) {ecDHFR};
\end{axis}

% =====================================================================
% (B) the offset against temperature
% =====================================================================
\begin{axis}[env, name=B, width=31mm, height=50mm,
  at={($(A.east)+(16mm,0)$)}, anchor=west,
  xmin=250, xmax=350, ymin=-0.04, ymax=1.12,
  xlabel={temperature (K)},
  ylabel={$[F(T)-F(250\,\mathrm K)]/|\Fz|$},
  xtick={250,300,350}, ytick={0,0.25,0.5,0.75,1},
  yticklabels={$0$,$0.25$,$0.5$,$0.75$,$1$},
  title={(\textit{\textbf{B}})\ \ gating moves $F$\\
         \hspace*{1.2mm}with temperature}]
  \addplot[cGrey, line width=0.9pt] table[x=T,y=dF] {data/fe_T_50.dat};
  \addplot[cHyd, line width=0.9pt] table[x=T,y=dF] {data/fe_T_200.dat};
  \addplot[cPro, line width=0.9pt] table[x=T,y=dF] {data/fe_T_400.dat};
  \addplot[cGold, line width=0.9pt] table[x=T,y=dF] {data/fe_T_800.dat};
  \node[ann, anchor=south east] at (axis cs:347,1.005) {$400$};
  \node[ann, anchor=south east, color=cHyd] at (axis cs:347,0.855) {$200$};
  \node[ann, anchor=south east, color=cGold] at (axis cs:347,0.355) {$800$};
  \node[ann, anchor=south east, color=cGrey] at (axis cs:347,0.080) {$50$};
  \node[ann, anchor=north west] at (axis cs:254,1.07)
    {mode (cm$^{-1}$)\\$\mathcal{A}=5$, $w_0=0.5$};
\end{axis}

% =====================================================================
% (C) where the readout is sensitive
% =====================================================================
\begin{axis}[env, name=C, width=37mm, height=50mm,
  at={($(B.east)+(15mm,0)$)}, anchor=west,
  xmode=log, xmin=30, xmax=1200, ymin=0, ymax=1.32,
  xlabel={promoting-mode wavenumber $\tilde\nu_g$ (cm$^{-1}$)},
  ylabel={span of $F$, $250$--$350$~K, in $|\Fz|$},
  xtick={30,100,300,1000}, xticklabels={$30$,$100$,$300$,$1000$},
  ytick={0,0.25,0.5,0.75,1},
  yticklabels={$0$,$0.25$,$0.5$,$0.75$,$1$},
  log basis x=10, clip=false,
  title={(\textit{\textbf{C}})\ \ most sensitive near\\
         \hspace*{1.2mm}the crossover}]
  % range of donor-acceptor frequencies invoked in the literature
  \fill[cGrey!14] (axis cs:50,0) rectangle (axis cs:300,1.32);
  \node[annc, color=cGrey!90!black] at (axis cs:122,0.07) {invoked\\$50$--$300$};
  \draw[cInk, line width=0.5pt, densely dashed]
    (axis cs:\feCross,0) -- (axis cs:\feCross,1.32);
  \node[ann, anchor=north west, rotate=90] at (axis cs:\feCross*1.07,0.08)
    {$hc\tilde\nu_g=2k_{\mathrm B}T$};
  \addplot[cPro, line width=1.0pt] table[x=nu,y=span] {data/fe_span.dat};
  \addplot[only marks, mark=*, mark size=1.5pt, cPro]
    table[x=nu,y=span] {data/fe_span_pts.dat};
  % where the peak sits over the sweep of A and w0
  \draw[cAcc, line width=0.8pt, |-|] (axis cs:\fePeakLo,1.21) -- (axis cs:\fePeakHi,1.21);
  \node[ann, color=cAcc!85!black, anchor=east, align=right] at (axis cs:130,1.21)
    {peak over\\$\mathcal{A}$, $w_0$};
\end{axis}

\end{tikzpicture}
\end{document}
